% !TeX program = xelatex
% =====================================================================
%  راهنمای دستور زبان لاتک
%  کدهای ارائه‌شده در جزوه‌ی «جزوه آموزشی دستور زبان لاتک»
%  (ویژه‌ی پشتیبانی زیپرشین / Zip Persian Support)
%  نویسنده‌ی جزوه: S. M. A. Khatami
%  کامپایل با XeLaTeX  (بسته xepersian)
%  نکته: در متن فارسی، عبارت‌های لاتین باید داخل \lr{...} نوشته شوند.
% =====================================================================
\documentclass[12pt,a4paper]{article}

\usepackage[margin=2.5cm]{geometry}
\usepackage{amsmath,amssymb,amsthm}
\usepackage{booktabs}
\usepackage{array}
\usepackage{enumitem}
\usepackage[table]{xcolor}
\usepackage{graphicx}
\usepackage{listings}
\usepackage[hidelinks]{hyperref}

% ---------- بسته فارسی (آخرین بسته) ----------
\usepackage{xepersian}
\settextfont[Scale=1.1,BoldFont={B Nazanin Bold}]{B Nazanin}
\setlatintextfont{Times New Roman}
\setmonofont{Courier New}

% نشانه‌های فهرست با فونت ریاضی (چون فونت B Nazanin گلیف بولت ندارد)
\renewcommand{\labelitemi}{\(\bullet\)}
\renewcommand{\labelitemii}{\(\bullet\)}
\renewcommand{\labelitemiii}{\(\bullet\)}

% ---------- تنظیمات نمایش کد ----------
\definecolor{codebg}{RGB}{247,247,247}
\definecolor{codekey}{RGB}{0,85,150}
\definecolor{codestr}{RGB}{160,80,0}
\definecolor{codecom}{RGB}{0,110,60}
\definecolor{codefr}{RGB}{190,190,190}

\lstdefinestyle{latexcode}{
  language=TeX,
  backgroundcolor=\color{codebg},
  basicstyle=\ttfamily\footnotesize,
  keywordstyle=\color{codekey}\bfseries,
  commentstyle=\color{codecom}\itshape,
  stringstyle=\color{codestr},
  numbers=left,
  numberstyle=\tiny\color{gray},
  numbersep=8pt,
  frame=single,
  rulecolor=\color{codefr},
  framerule=0.4pt,
  framesep=5pt,
  breaklines=true,
  breakatwhitespace=false,
  showstringspaces=false,
  tabsize=4,
  columns=fullflexible,
  keepspaces=true,
  upquote=true,
  aboveskip=10pt,
  belowskip=10pt,
}
\lstset{style=latexcode}

% ---------- راحتی‌ها ----------
\newcommand{\code}[1]{\lr{\texttt{\small #1}}}
\newcommand{\enp}[1]{\lr{(#1)}}
\renewcommand{\lstlistingname}{کد}

% محیط توضیحات (نکات) با نشان مربع توخالی
\newenvironment{tnotes}
  {\begin{itemize}\setlength{\itemsep}{3pt}\setlength{\parskip}{0pt}}
  {\end{itemize}}

% =====================================================================
\begin{document}

\begin{titlepage}
\centering
\vspace*{2cm}

{\Large دانشگاه ............ \par}
\vspace{0.5cm}
{\large دانشکده مهندسی کامپیوتر \par}
\vspace{2cm}

{\Huge\bfseries راهنمای دستور زبان لاتک \par}
\vspace{0.8cm}
{\LARGE\bfseries آموزش نگارش متن فارسی، فرمول و جدول \par}
\vspace{0.4cm}
{\large\lr{How to type with Latex} \par}

\vspace{2cm}
\begin{tabular}{r r}
\textbf{نویسنده‌ی جزوه} & \lr{S. M. A. Khatami} \\
\textbf{تاریخ ارائه} & آبان ۱۴۰۵ \\
\textbf{منبع} & \lr{Zip Persian Support} \\
\end{tabular}

\vfill
\end{titlepage}

\tableofcontents
\newpage

% =====================================================================
\section*{مقدمه و تنظیمات اولیه}
\addcontentsline{toc}{section}{مقدمه و تنظیمات اولیه}

در این جزوه، ساختار استاندارد یک سند \lr{LaTeX} با پشتیبانی کامل از زبان فارسی
(بسته‌ی \lr{xepersian}) معرفی می‌شود. سپس ساختار لیست‌ها، نوشتن فرمول‌های ریاضی،
آدیه‌ها و ماتریس‌ها و در پایان جدول‌های ساده و رنگی آموزش داده می‌شود.
همه‌ی کدهای این جزوه کامل و بدون خطا هستند و می‌توان آن‌ها را مستقیماً در
\lr{TeXstudio} کپی و کامپایل کرد.

% =====================================================================
\section{فرم‌های پکیج‌ها و ساختار سند}

کد زیر، اسکلت استاندارد یک سند فارسی با \lr{xepersian} است. دقت کنید که بسته‌ی
فارسی باید آخرین بسته‌ای باشد که بارگذاری می‌شود.

\begin{latin}
\begin{lstlisting}[caption={اسکلت استاندارد سند فارسی}]
\documentclass{article}
\usepackage{xepersian}
\settextfont{Yas}
\usepackage{lipsum,amssymb,amsfonts}
\usepackage{enumitem}
\usepackage[table]{xcolor}
\usepackage{array}

\title{How to type with Latex}
\author{S. M. A. Khatami}
\date{Mehr 1405}

\begin{document}
    \maketitle
\end{document}
\end{lstlisting}
\end{latin}

\begin{tnotes}
\item \textbf{نکته‌ی ۱:} بسته‌ی \lr{xepersian} برای پشتیبانی از زبان فارسی است و
      باید پس از همه‌ی بسته‌های دیگر بارگذاری شود.
\item \textbf{نکته‌ی ۲:} دستور \lr{\textbackslash settextfont\{Yas\}} نام فونت متن را
      تعیین می‌کند؛ می‌توانید نام فونت دیگری را که روی سیستم خود نصب است بنویسید
      (در این پروژه از \lr{B Nazanin} استفاده شده است).
\item \textbf{نکته‌ی ۳:} بسته‌های کمکی برای ریاضیات \lr{amssymb} و \lr{amsfonts}،
      لیست‌ها \lr{enumitem}، رنگ‌آمیزی جدول \lr{xcolor} با گزینه‌ی \lr{table} و
      ساختار ستون‌های جدول \lr{array} آماده شده‌اند.
\item \textbf{نکته‌ی ۴:} دستورات \lr{\textbackslash author}، \lr{\textbackslash title} و
      \lr{\textbackslash date} اطلاعات عنوان، نویسنده و تاریخ را ذخیره می‌کنند که با
      دستور \lr{\textbackslash maketitle} در صفحه‌ی اول چاپ می‌شوند.
\end{tnotes}

% =====================================================================
\section{لیست‌ها و شماره‌گذاری}

گاهی لازم است شماره‌گذاری یک فهرست از عددی غیر از ۱ آغاز شود. بسته‌ی
\lr{enumitem} این امکان را با گزینه‌ی \lr{start} فراهم می‌کند.

\begin{latin}
\begin{lstlisting}[caption={فهرست شماره‌دار با شروع از ۱۲ و لیست تودرتو}]
\begin{enumerate}[start=12]
    \item download texlive
    \item
    \begin{itemize}
        \item install a virtual cd drive
        \item mount iso image of texlive
        \item install texlive
    \end{itemize}
    \item download and install TexStdio
\end{enumerate}
\end{lstlisting}
\end{latin}

خروجی این کد به صورت زیر است:

\begin{latin}
\begin{enumerate}[start=12]
    \item download texlive
    \item
    \begin{itemize}
        \item install a virtual cd drive
        \item mount iso image of texlive
        \item install texlive
    \end{itemize}
    \item download and install TexStdio
\end{enumerate}
\end{latin}

\begin{tnotes}
\item \textbf{نکته‌ی ۱:} گزینه‌ی \lr{[start=12]} باعث می‌شود شماره‌گذاری فهرست از عدد
      ۱۲ شروع شود.
\item \textbf{نکته‌ی ۲:} یک لیست تودرتو (\lr{itemize}) درون محیط \lr{enumerate} ایجاد
      شده است که آیتم‌های آن با نشان بولت نمایش داده می‌شوند.
\end{tnotes}

% =====================================================================
\section{فرمول‌های ریاضی}

\subsection{حل معادله‌ی درجه دوم}

برای حل معادله‌ی $x^2-5x+6=0$ دو روش می‌توان به‌کار برد: روش اول بر پایه‌ی
equalities جبری و روش دوم بر پایه‌ی فرمول دلتا. کد زیر هر دو روش را نشان می‌دهد.

\begin{latin}
\begin{lstlisting}[caption={نوشتن فرمول ریاضی در متن}]
How solve the equatioin $x^2-5x+6=0$:
\begin{itemize}[leftmargin=2cm]
    \item[Method 1:] use \underline{algebraic} equalities
        $$x^2-5x+6=(x-2)(x-3) $$
        So $(x-2)(x-3) =0$ and \textbf{ therefore}
        $x=2$ or $x=3$.
    \item[Method 2:] Use the $\Delta $ method . For
        equation $ax ^2+ bx+c=0$ set $\Delta =b^2-4ac$ and
        then the roots are $$x=\frac{-b\mp\sqrt{ \Delta }}{2a}$$ here $\Delta =(-5) ^2-4(1) (6) =25-24=1$ and so
        $$x=\frac{-(-5)\mp\sqrt{1}}{2(1) }=\frac{5\mp 1}{2}$$
    \[
    \begin{array}{c}
        \frac{5+1}{2}=3\\
        \frac{5- 1}{2}=2
    \end{array}
    \]
\end{itemize}
\end{lstlisting}
\end{latin}

خروجی این کد:

\begin{latin}
\begin{itemize}[leftmargin=2cm]
    \item[Method 1:] use \underline{algebraic} equalities
    \[
        x^2-5x+6=(x-2)(x-3)
    \]
    So $(x-2)(x-3) =0$ and \textbf{ therefore}
    $x=2$ or $x=3$.
    \item[Method 2:] Use the $\Delta $ method . For
    equation $ax ^2+ bx+c=0$ set $\Delta =b^2-4ac$ and
    then the roots are
    \[
        x=\frac{-b\mp\sqrt{ \Delta }}{2a}
    \]
    here $\Delta =(-5) ^2-4(1) (6) =25-24=1$ and so
    \[
        x=\frac{-(-5)\mp\sqrt{1}}{2(1) }=\frac{5\mp 1}{2}
    \]
    \[
    \begin{array}{c}
        \frac{5+1}{2}=3\\
        \frac{5- 1}{2}=2
    \end{array}
    \]
\end{itemize}
\end{latin}

\begin{tnotes}
\item \textbf{نکته‌ی ۱:} گزینه‌ی \lr{[leftmargin=2cm]} حاشیه‌ی چپ لیست را به اندازه‌ی
      ۲ سانتی‌متر تنظیم می‌کند.
\item \textbf{نکته‌ی ۲:} علامت \lr{\$} برای فرمول‌های درون‌خطی و \lr{\$\$} برای
      فرمول‌های نمایشی (وسط‌چین) استفاده می‌شود.
\item \textbf{نکته‌ی ۳:} دستور \lr{\textbackslash underline{}} زیر متن را خط می‌کشد و
      دستور \lr{\textbackslash textbf{}} متن را پررنگ می‌کند.
\item \textbf{نکته‌ی ۴:} برای نوشتن معادلات در چند سطر از محیط \lr{array} درون
      \lr{\textbackslash left\{ ... \textbackslash right.} استفاده شده است.
      \lr{\textbackslash left.} و \lr{\textbackslash right} اندازه‌ی اکولادهای
      را با محتوای داخل محیط تطبیق می‌دهند.
\item \textbf{نکته‌ی ۵:} دستور \lr{\textbackslash frac} برای کسر و \lr{\textbackslash sqrt}
      برای ریشه‌ی دوم به‌کار می‌رود.
\end{tnotes}

\subsection{آدیه‌ها و ماتریس‌ها}

برای رسم ماتریس‌ها از محیط \lr{array} استفاده می‌کنیم:

\begin{latin}
\begin{lstlisting}[caption={ساخت ماتریس با محیط array}]
\[
\begin{array}{ cc}
    1 & 2 \\
    3 & 8 \\
    2 & x^2+x-1+ \frac{1}{4}
\end{array}
\]
\end{lstlisting}
\end{latin}

خروجی این کد:

\begin{latin}
\[
\begin{array}{ cc}
    1 & 2 \\
    3 & 8 \\
    2 & x^2+x-1+ \frac{1}{4}
\end{array}
\]
\end{latin}

\begin{tnotes}
\item \textbf{نکته‌ی ۱:} دستور \lr{\textbackslash left} یک براکت بزرگ در سمت چپ محیط
      ایجاد می‌کند.
\item \textbf{نکته‌ی ۲:} در \lr{\{cc\}} یعنی دو ستون و دو سطر. علامت \lr{\&} ستون‌ها
      را از هم جدا می‌کند و \lr{\textbackslash\textbackslash} سطر جدید ایجاد می‌کند.
\item \textbf{نکته‌ی ۳:} در نسخه‌ی اصلی جزوه، دستور \lr{\textbackslash end\{array\}} در
      انتهای سطر آخر قرار گرفته بود که اصلاح شد تا کد بدون خطا کامپایل شود.
\end{tnotes}

% =====================================================================
\section{جداول پیشرفته}

\subsection{جدول ساده با ادغام ستون‌ها}

\begin{latin}
\begin{lstlisting}[caption={جدول ساده با خط‌کشی و ادغام ستون}]
\begin{center}
    \begin{tabular}{|c||c|c|c|c|c|c|}
        \hline
        &8-10&10 - 12&12 -2&2-4&4-6&6-8\\
        \hline\hline
        sat & & & & & & Spesial topics \\
        \hline
        sun & Spesial topics & & & && \\
        \hline
    \end{tabular}
\end{center}
\end{lstlisting}
\end{latin}

خروجی این کد:

\begin{latin}
\begin{center}
    \begin{tabular}{|c||c|c|c|c|c|c|}
        \hline
        &8-10&10 - 12&12 -2&2-4&4-6&6-8\\
        \hline\hline
        sat & & & & & & Spesial topics \\
        \hline
        sun & Spesial topics & & & && \\
        \hline
    \end{tabular}
\end{center}
\end{latin}

\begin{tnotes}
\item \textbf{نکته‌ی ۱:} الگوی \lr{|c|c|...} یعنی ستون‌ها و سطرها هستند (خط‌کشیده
      می‌شوند) و \lr{||} یعنی خط دوبل.
\item \textbf{نکته‌ی ۲:} دستور \lr{\textbackslash hline} خط افقی می‌کشد. در پایان
      سطر هم \lr{\textbackslash hline\textbackslash hline} خط دوبل می‌کشد.
\end{tnotes}

\subsection{جدول رنگی با عرض ثابت}

\begin{latin}
\begin{lstlisting}[caption={جدول رنگی با عرض ستون ثابت}]
\begin{center}
    \begin{tabular}{|c ||*{6}{>{\centering\arraybackslash}p{1.8cm}|}}
        \hline
        \rowcolor{gray!30}
        & 8-10 & 10-12 & 12-2 & 2-4 & 4-6 & 6-8 \\
        \hline\hline

        \cellcolor{gray!30} sat
        & & & & & &
        Special topics \\
        \hline

        \cellcolor{gray!30} sun
        & Special topics & & & & &
        Special topics\\
        \cline{7-7}
        \cellcolor{gray!30}
        & & & & & &
        Computability \\
        \hline
    \end{tabular}
\end{center}
\end{lstlisting}
\end{latin}

خروجی این کد:

\begin{latin}
\begin{center}
    \begin{tabular}{|c ||*{6}{>{\centering\arraybackslash}p{1.8cm}|}}
        \hline
        \rowcolor{gray!30}
        & 8-10 & 10-12 & 12-2 & 2-4 & 4-6 & 6-8 \\
        \hline\hline

        \cellcolor{gray!30} sat
        & & & & & &
        Special topics \\
        \hline

        \cellcolor{gray!30} sun
        & Special topics & & & & &
        Special topics\\
        \cline{7-7}
        \cellcolor{gray!30}
        & & & & & &
        Computability \\
        \hline
    \end{tabular}
\end{center}
\end{latin}

\begin{tnotes}
\item \textbf{نکته‌ی ۱:} الگوی \lr{*{6}{...}} یعنی ۶ ستون با تنظیماتی که در
      \lr{\{>{\textbackslash centering\textbackslash arraybackslash\}p\{1.8cm\}\}}
      مشخص شده است؛ یعنی یک پارگراف با عرض ۱.۸ سانتی‌متر که متن آن وسط‌چین است.
\item \textbf{نکته‌ی ۲:} دستور \lr{\textbackslash rowcolor\{gray!30\}} کل سطر را با
      رنگ خاکستری ۳۰ درصد رنگ‌آمیزی می‌کند.
\item \textbf{نکته‌ی ۳:} دستور \lr{\textbackslash cellcolor\{gray!30\}} فقط یک خانه را
      رنگ می‌کند.
\item \textbf{نکته‌ی ۴:} دستور \lr{\textbackslash cline\{7-7\}} یک خط افقی فقط در ستون
      هفتم ایجاد می‌کند تا دو آیتم داخل یک خانه از هم جدا شوند.
\end{tnotes}

% =====================================================================
\section{جمع‌بندی}

\begin{itemize}
\item بسته‌ی \lr{xepersian} باید آخرین بسته باشد و فونت متن با
      \lr{\textbackslash settextfont} تعیین می‌شود.
\item بسته‌ی \lr{enumitem} کنترل کاملی روی شماره‌گذاری، فاصله‌ها و حاشیه‌های
      لیست‌ها می‌دهد.
\item فرمول‌ها با \lr{\$} (درون‌خطی) و \lr{\$\$} یا \lr{\textbackslash[ ... \textbackslash]}
      (نمایشی) نوشته می‌شوند.
\item محیط \lr{array} برای ماتریس‌ها و محیط \lr{tabular} برای جدول‌ها به‌کار می‌رود؛
      با بسته‌ی \lr{xcolor} و گزینه‌ی \lr{table} می‌توان سطرها و خانه‌ها را رنگی کرد.
\end{itemize}

\end{document}
